% chapters/abstract.tex
\addcontentsline{toc}{section}{АНОТАЦІЯ}
{\centering\bfseries АНОТАЦІЯ\par}
\vspace{1em}

\textbf{Тесля М.} Візуальне моделювання роботи обчислювача ШПФ під
керуванням потоку даних. Бакалаврська кваліфікаційна робота за
спеціальністю 122 «Комп'ютерні науки». Національний університет
«Львівська політехніка», Львів, 2026.

У межах бакалаврської кваліфікаційної роботи розроблено математичну
модель, граф потоку даних та веб-орієнтований програмний симулятор
16-точкового швидкого перетворення Фур'є за алгоритмом radix-4 для
інтерактивної візуалізації потокового виконання обчислень із
використанням технологій Next.js, React~Flow, Zustand і TypeScript.

Робота складається із вступу, чотирьох основних розділів, економічної
частини, висновків, списку використаних джерел та додатків.

У вступі обґрунтовано актуальність теми, визначено мету, задачі,
об'єкт і предмет дослідження, описано методологію та інструментарій,
а також практичну цінність розробки.

У першому розділі проведено аналіз предметної області, огляд
сучасних підходів до обчислення ШПФ, порівняльний аналіз існуючих
рішень (FFTW, CMSIS-DSP, cuFFT, SPIRAL), виконано системний аналіз
із виділенням зацікавлених сторін, цілей та обмежень, розглянуто
етичні й правові аспекти та сформульовано формальну постановку
задачі з вимогами F1--F9, NF1--NF8 і сценаріями використання UC1--UC9.

Другий розділ присвячено проєктуванню системи: обґрунтовано вибір
клієнтського SPA-підходу, розроблено математичну модель тристадійної
$4\times4$ декомпозиції, сформовано граф потоку даних (56~вузлів,
64~дуги), визначено критичний шлях і коефіцієнт паралелізму, описано
структури даних та механізми обміну між компонентами.

Третій розділ описує програмну реалізацію симулятора: технологічний
стек, структуру проєкту, реалізацію ядра \texttt{DataflowRuntime}
з методом \texttt{tick()}, обчислювальних вузлів, алгоритму
\texttt{generateDFT4x4()}, класифікацію винятків, процедуру
розгортання та модульне тестування у середовищі Jest.

У четвертому розділі подано результати експериментальних досліджень:
верифікацію на шести тестових сценаріях (максимальна похибка
$3\times10^{-12}$ при допуску $\varepsilon=10^{-9}$), аналіз
обчислювальної складності ($K_M\approx28{,}4\times$), аналіз
паралелізму dataflow-моделі ($K_p\approx1{,}88\times$),
обговорення, моніторинг експлуатаційних метрик та підтвердження
нефункціональних вимог NF1--NF8.

У висновках узагальнено основні результати виконаної роботи,
визначено практичне значення та перспективні напрями подальших
досліджень.

Дана бакалаврська кваліфікаційна робота містить: \pageref{LastPage}~сторінок,
1~додаток та 25~джерел літератури.

\textbf{Ключові слова:} швидке перетворення Фур'є, radix-4, граф
потоку даних, dataflow computing, візуальне моделювання,
веб-симулятор, Next.js, React~Flow, паралелізм, цифрове оброблення
сигналів.

\vspace{1em}\hrule\vspace{1em}

\addcontentsline{toc}{section}{ABSTRACT}
{\centering\bfseries ABSTRACT\par}
\vspace{1em}

\begin{otherlanguage}{english}
\textbf{Teslia M.} Visual Modeling of a Dataflow-Driven FFT Processor.
Bachelor's Qualification Work for specialty 122 «Computer Science».
Lviv Polytechnic National University, Lviv, 2026.

Within this bachelor's qualification work, a mathematical model, a
data-flow graph, and a web-based software simulator of a 16-point
Fast Fourier Transform (radix-4) have been developed for interactive
visualization of the dataflow-driven computation, using Next.js,
React~Flow, Zustand, and TypeScript technologies.

The work consists of an introduction, four main chapters, an economic
part, conclusions, a list of references, and appendices.

The introduction substantiates the relevance of the topic and defines
the aim, tasks, object, subject of the study, the methodology, and
the practical value of the development.

The first chapter presents an analysis of the problem domain, a
review of modern FFT computation approaches, a comparative analysis
of existing solutions (FFTW, CMSIS-DSP, cuFFT, SPIRAL), a systems
analysis with stakeholders, goals and constraints, a discussion of
ethical and legal aspects, and a formal task statement with F1--F9,
NF1--NF8 requirements and UC1--UC9 use cases.

The second chapter is devoted to the system design: justification of
a client-side SPA approach, a mathematical model of the three-stage
$4\times4$ decomposition, a data-flow graph (56~nodes, 64~arcs),
determination of the critical path and parallelism factor, data
structures, and inter-component communication mechanisms.

The third chapter describes the software implementation of the
simulator: technology stack, project structure, the
\texttt{DataflowRuntime} engine with the \texttt{tick()} method,
computational nodes, the \texttt{generateDFT4x4()} algorithm,
exception handling, deployment procedure, and unit testing using Jest.

The fourth chapter presents the experimental results: verification
across six test scenarios (maximum error $3\times10^{-12}$ against
tolerance $\varepsilon=10^{-9}$), analysis of computational complexity
($K_M\approx28.4\times$), analysis of dataflow-model parallelism
($K_p\approx1.88\times$), discussion, monitoring of operational
metrics, and confirmation of non-functional requirements NF1--NF8.

The conclusions summarize the key results of the work and outline
the practical value and promising directions for further research.

This bachelor's qualification work contains: \pageref{LastPage}~pages,
1~appendix, and 25~references.

\textbf{Keywords:} fast Fourier transform, radix-4, data-flow graph,
dataflow computing, visual modeling, web-based simulator, Next.js,
React~Flow, parallelism, digital signal processing.
\end{otherlanguage}
